\section{Metrics}
\label{sec:metrics}

\subsection{Notation}

Fix a population, a sex and an origin $T$. For age $x$ and horizon $h$ the
target on the rate scale is $y_{x,h} = \log m_{x,T+h}$ and on the count scale
it is the observed death count $D_{x,T+h}$ with known exposure $E_{x,T+h}$.
A cell of the grid supplies $n$ predictive sample paths
$z^{(j)}_{x,h}=\log m^{(j)}_{x,T+h}$, $j=1,\dots,n$, jointly over $(x,h)$;
$\hat F_{x,h}$ denotes the empirical predictive distribution of
$\{z^{(j)}_{x,h}\}_j$ and $\hat F^{-1}_{x,h}(p)$ its empirical quantile.
All scores are negatively oriented: smaller is better. Averages over cells
are denoted by an overbar; averaging is always within population--sex
first, then across populations, so that a large population cannot dominate
a small one and so that the cluster structure of Section~\ref{sec:metrics-inference}
is preserved.

\subsection{Point accuracy}

\begin{equation}
  \RMSE = \Big(\overline{(\hat y_{x,h}-y_{x,h})^2}\Big)^{1/2},\qquad
  \MAE = \overline{|\hat y_{x,h}-y_{x,h}|},\qquad
  \hat y_{x,h} = \hat F^{-1}_{x,h}(\tfrac12),
\end{equation}
together with the error in period life expectancy at birth,
$\hat e_0 - e_0$, computed from the median path. These are the metrics the
literature reports; \Hone\ is the claim that ranking by them is not ranking
by a proper score.

\subsection{Continuous ranked probability score}

For a predictive distribution $F$ and outcome $y$,
\begin{equation}
  \CRPS(F,y) = \int_{-\infty}^{\infty}\big(F(u)-\ind\{u\ge y\}\big)^2\,du
  = \E_F|Z-y| - \tfrac12\,\E_F|Z-Z'|,
\end{equation}
with $Z,Z'\sim F$ independent \citep{gneiting2007strictly}. From samples we
use the plug-in estimator with sorted samples $z_{(1)}\le\dots\le z_{(n)}$,
\begin{equation}
  \widehat{\CRPS} = \frac1n\sum_{j=1}^{n}|z^{(j)}-y|
  - \frac{1}{n^2}\sum_{j=1}^{n}\big(2j-n-1\big)\,z_{(j)},
  \label{eq:crps-sample}
\end{equation}
which is the $O(n\log n)$ form of the second term. CRPS is computed on the
log-rate scale, $y=y_{x,h}$, for every cell, and on the count scale as a
registered sensitivity (Section~\ref{sec:metrics-logscore}). It is strictly
proper and reduces to MAE for a point forecast, which is what makes the
RMSE-versus-CRPS ranking comparison of \Hone\ well posed.

\subsection{Poisson log score and the rounding convention}
\label{sec:metrics-logscore}

Each sample path implies a Poisson mean $\lambda^{(j)}_{x,h} =
E_{x,T+h}\,m^{(j)}_{x,T+h}$ for the death count, so the predictive count
distribution is the equal-weight mixture
$p(d)=\frac1n\sum_j \mathrm{Poisson}(d;\lambda^{(j)}_{x,h})$. The log score is
\begin{equation}
  \LS = -\log\Big(\frac1n\sum_{j=1}^{n}\mathrm{Poisson}\big(\tilde D_{x,T+h};\lambda^{(j)}_{x,h}\big)\Big),
  \qquad \tilde D = \big\lfloor D+\tfrac12\big\rfloor,
\end{equation}
evaluated by log-sum-exp. Because HMD deaths are fractional
(Section~\ref{sec:data}), the score is defined at the nearest integer;
the sensitivity is the CRPS of the mixture count distribution against the
unrounded $D$, which is well defined for any real $D$. Both conventions were
fixed before estimation.

\subsection{Interval score and coverage}

For the central $(1-\alpha)$ interval $[l,u] =
[\hat F^{-1}(\alpha/2),\hat F^{-1}(1-\alpha/2)]$ the interval (Winkler) score
is
\begin{equation}
  \IS_\alpha(l,u;y) = (u-l) + \frac{2}{\alpha}(l-y)\,\ind\{y<l\}
  + \frac{2}{\alpha}(y-u)\,\ind\{y>u\},
  \label{eq:winkler}
\end{equation}
\citep{winkler1972decision,gneiting2007strictly}, reported at $\alpha\in\{0.5,0.2,0.05\}$.
It is the width-penalised criterion whose absence from the mortality
literature Section~\ref{sec:intro} documents: an interval wide enough to
cover everything has coverage one and an interval score equal to its width.
Alongside it we always report the pair (empirical coverage, mean width),
never coverage alone.

\subsection{Marginal and joint path coverage}
\label{sec:metrics-coverage}

\emph{Marginal} coverage is the pointwise hit rate,
\begin{equation}
  C^{\mathrm{marg}}_{1-\alpha}(h) = \overline{\ind\{\,l_{x,h}\le y_{x,h}\le u_{x,h}\,\}},
\end{equation}
averaged over ages and populations for each horizon $h$, and pooled over
$h$ for a single number. It is the quantity every mortality paper that
reports coverage reports.

\emph{Joint path} coverage asks whether the whole trajectory of a unit
(one population--sex--age) lies inside its bands:
\begin{equation}
  C^{\mathrm{joint}}_{1-\alpha} = \overline{\prod_{h=1}^{H}\ind\{\,l_{x,h}\le y_{x,h}\le u_{x,h}\,\}},
  \qquad H=5 \text{ (shift, stable)},\ H=9 \text{ (placebo)}.
\end{equation}
A path counts as covered only if every horizon is inside its band. This is
the object an annuity valuation actually depends on, because a liability is
a functional of the path of future mortality, not of one year of it. For
pointwise bands at level $1-\alpha$, joint coverage is bounded above by
marginal coverage and, if horizons were independent, would equal
$(1-\alpha)^H$; the strong horizon dependence of random-walk paths puts it
between the two. Only the copula conformal mechanism targets the joint level
by construction. \Hthree\ is the registered claim that the gap
$C^{\mathrm{marg}}-C^{\mathrm{joint}}$ is material for every family, tested
as a paired difference over units.

\subsection{Randomised probability integral transform}

A forecaster is probabilistically calibrated if $F(Y)$ is uniform
\citep{gneiting2007probabilistic,dawid1984present}. Because the predictive is
an empirical distribution of $n$ samples and the target may tie with them,
we use the randomised PIT of \citet{czado2009predictive},
\begin{equation}
  u_{x,h} = \frac{\#\{j: z^{(j)}_{x,h} < y_{x,h}\} + V\,\big(\#\{j: z^{(j)}_{x,h}=y_{x,h}\}+1\big)}{n+1},
  \qquad V\sim\mathrm{Uniform}(0,1),
\end{equation}
which is exactly uniform under the null for any finite $n$. We report the
ten-bin histogram, the Kolmogorov--Smirnov distance from uniformity, and the
shape---U (over-confident), hump (under-confident), slope (biased). Because
PIT values within a population--sex--origin are strongly dependent, the
formal test is a cluster-bootstrap test of the pooled histogram
(Section~\ref{sec:metrics-inference}), not a cell-level KS $p$-value.

\subsection{Calibration by age}

All coverage and PIT statistics are additionally computed within the age
bands $0$--$24$, $25$--$64$ and $65$--$99$ and as single-age curves.
\Hfour\ compares the $65$--$99$ and $25$--$64$ bands in the shift regime.

\subsection{Murphy decomposition}

For a proper score $S$ the mean score decomposes as
\begin{equation}
  \overline{S} = \mathrm{UNC} - \mathrm{RES} + \mathrm{MCB},
\end{equation}
uncertainty of the outcome minus resolution of the forecasts plus
miscalibration \citep{murphy1973new,pohle2020murphy}. It separates
\emph{why} a forecaster scores badly: a well-calibrated but vague
forecaster has low RES; a sharp but wrong one has high MCB. We compute it
for the interval score and for the Brier scores of the exceedance events
$\ind\{Y\le \hat F^{-1}(p)\}$ at $p\in\{0.025,0.25,0.5,0.75,0.975\}$, with
the recalibrated forecast estimated by isotonic regression
\citep{dimitriadis2021stable}. \todo{estimator not yet implemented in the
harness; confirm the isotonic (CORP) choice or a binned alternative before
the first real-data run}

\subsection{Actuarial functionals}
\label{sec:metrics-actuarial}

Each sample path's rate vector $m^{(j)}_{\cdot,T+h}$ is converted to a period
life table by the conventions of Section~\ref{sec:data}, giving samples of
\begin{equation}
  e^{(j)}_x = \frac{T^{(j)}_x}{l^{(j)}_x},\qquad
  \annuity^{(j)} = \sum_{t\ge0} v^{t}\,{}_t p^{(j)}_{65},\quad v=\frac{1}{1.02},
\end{equation}
where ${}_t p_{65}=l_{65+t}/l_{65}$ from the sample's own table treated as
static, i.e.\ the annuity-due factor at 2\% on the period table of year
$T+h$ \citep{dickson2020actuarial}. The realised values are computed from the
observed $m_{x,T+h}$ by the same routine, so that any life-table
approximation cancels in the comparison. For $e_0$, $e_{65}$ and $\annuity$
we report the interval error, the empirical coverage of the central 95\%
interval and its width, in each regime; \Hfive\ compares the shift and
stable departures from nominal. We also report the 99.5\% upper quantile of
$\annuity$ from each cell as the longevity tail an internal model would
carry. \todo{the 99.5\% tail quantile is descriptive and not part of a
registered hypothesis; say so in the table note}

\subsection{Forecast-comparison inference}
\label{sec:metrics-inference}

\paragraph{Effective sample size.} The shift regime contains one common
shock. Within a population, ages and horizons share the same realised
pandemic; across populations, the shock itself is shared. Treating cells as
independent would produce standard errors that are wrong by orders of
magnitude. Population is therefore the cluster ($G=20$ in the shift and
stable regimes, $G=11$ in the placebo), and every test resamples clusters.
Even so, $p$-values in the shift regime are statements about differences
between forecasters \emph{given this event}, not about future pandemics;
that is what an event study can deliver, and the placebo regime is the
second event that makes a cross-event comparison possible.

\paragraph{Diebold--Mariano with a wild cluster bootstrap.} For forecasters
$A$ and $B$ with per-unit losses $L_{A,i}$, $L_{B,i}$ (CRPS per horizon per
population--sex--age), the differential $d_i = L_{A,i}-L_{B,i}$ has mean
$\bar d$ and cluster-robust variance
\begin{equation}
  \widehat{\mathrm{Var}}(\bar d) = \frac{G}{G-1}\,\frac{1}{N^2}\sum_{g=1}^{G}
  \Big(\sum_{i\in g}(d_i-\bar d)\Big)^2,
\end{equation}
giving the statistic $t=\bar d/\widehat{\mathrm{se}}$
\citep{diebold1995comparing}. Its null distribution is obtained by the wild
cluster bootstrap of \citet{cameron2008bootstrap} with the null imposed:
residuals $d_i-\bar d$ are multiplied by a cluster-level weight $w_g$ drawn
from the six-point distribution of \citet{webb2023reworking}, which has
enough distinct sign patterns at $G=20$ where Rademacher weights do not
\citep{mackinnon2017wild}; $B=4999$ replications. Negative $\bar d$ favours
$A$.

\paragraph{Model confidence set.} Following \citet{hansen2011model}, the MCS
at confidence $1-\alpha=0.90$ is built by the $T_{\max}$ statistic and the
$e_{\max}$ elimination rule, with the bootstrap distribution of the
statistic obtained by resampling \emph{populations} as blocks---the analogue
of their stationary block bootstrap for our panel. The MCS is reported for
CRPS on log rates in each regime, over the distributional cells of the grid.

\paragraph{Hypothesis tests.} \Hone\ is assessed by the Spearman correlation
of RMSE and CRPS rankings and the count of top-3 inversions, with a
cluster-bootstrap interval for the correlation. \Htwo, \Hfour\ and \Hfive\
are paired contrasts of coverage (shift minus stable; ages 65--99 minus
25--64) tested by the wild cluster bootstrap above applied to the coverage
indicators. \Hthree\ is the paired difference of marginal and joint
coverage over units, clustered by population.
